\section{GlyphNet Sovereign Communications and GlyphChain Proof Layer}
\label{sec:pilot-glyphnet-proof-layer}

\subsection{Decision and Product Boundary}

Pilot will reuse the substantial communications, identity, transport and ledger
infrastructure already present in GlyphNet, GlyphChain and the wider AION
runtime. It will not create a second, incompatible messaging stack. GlyphNet
becomes the common communications substrate for Personal Pilot, external
workspaces and Boardroom, while GlyphChain becomes a privacy-preserving proof
rail for consequential requests and actions.

The first product release will expose familiar capabilities---conversations,
voice notes, push-to-talk, tasks, approvals and agent-to-agent requests---through
the unified Pilot application. The implementation may use GlyphNet internally
without requiring an ordinary user to understand Wave addresses, radio routing,
blockchain terminology or token economics.

The PHO currency, PhotonPay, TESS, wrapped GLYPH, Glyph bonds, experimental
wallet balances and related economic controls will remain dormant behind
separate build-time and run-time feature flags. Their code and schemas are to be
preserved for future evaluation, but they are not dependencies of messaging,
identity, task delivery, service execution or the Pilot adoption loop.

\begin{quote}
\textbf{Decided architecture:} GlyphNet carries human and agent communication.
GlyphChain proves governed authorization and outcome. Existing payment
providers move money. PHO remains hidden until its security, economics,
settlement, custody and regulatory model are independently ready.
\end{quote}

\subsection{Existing Infrastructure to Retain}

The earlier implementation contains more than a visual prototype. The following
capabilities are sufficiently developed to form the starting point of the new
build, although each remains subject to the production hardening requirements
in this section:

\begin{itemize}
  \item typed one-to-one and threaded conversations;
  \item voice notes and browser microphone capture;
  \item press-and-hold push-to-talk with release-to-send behaviour;
  \item a floor-lock protocol that prevents simultaneous push-to-talk speakers,
        with bounded expiry and renewal;
  \item file-attachment, recipient and contact interface foundations;
  \item real-time WebSocket delivery, reconnect handling and delivery telemetry;
  \item WebRTC call signalling for offers, answers, network candidates,
        rejection, cancellation and termination;
  \item persistent local conversation and event journalling in SQLite;
  \item distinct personal and work conversation graphs;
  \item stable person or agent attribution separated from temporary browser-tab
        and transport identities;
  \item Wave-address and Wave-number identity concepts intended to span
        messaging, push-to-talk, devices and future payments;
  \item automatic selection between ordinary IP transport and an available
        radio bridge;
  \item bounded radio frames, fragmentation, queues, reconnection, local
        broadcast and store-and-forward foundations;
  \item offline mesh transaction and reconciliation data structures;
  \item wallet, PhotonPay and chain-simulation components retained as dormant
        experimental infrastructure; and
  \item persistent GlyphChain simulation and AION proof-commit and proof-lookup
        foundations.
\end{itemize}

This inventory records existing foundations, not production certification. In
particular, some development routes use fixed identities, placeholder
signatures, demonstration balances or mock transports. Those paths must never
be enabled merely because their user interface exists.

\subsection{One Communication Model Across Pilot}

Every user-visible communication will use one normalized envelope regardless of
the space in which it originates. The same envelope can represent a human
message, an agent request, a delegated task, an approval proposal, a reminder,
a call signal, a voice note or a provider handoff.

\begin{verbatim}
Pilot Inbox
  +-- Personal conversations
  +-- Household and trusted-contact channels
  +-- Pilot-to-Pilot requests
  +-- Tasks, reminders and approvals
  +-- External workspaces
  +-- Boardroom communications
            |
            v
     GlyphNet message envelope
            |
     +------+------+----------------+
     |             |                 |
     v             v                 v
 Local/VPN     Internet relay    Radio/mesh bridge
            |
            v
  Delivery or provider receipt
            |
            v
  Selective GlyphChain proof commitment
\end{verbatim}

The envelope will contain a version, opaque message identifier, conversation
identifier, sender identity, recipient identity or channel, originating space,
message type, creation and expiry times, reply and task relationships, transport
policy, encryption metadata, attachment references, idempotency key and a
content hash. Private content is encrypted payload, not ledger data.

Personal, household, workspace and Boardroom messages use the same protocol but
remain different authorization domains. Membership of one domain conveys no
right to read or act in another. A person's Pilot may participate in several
spaces, but the active space, role, capability lease and intended recipient must
be resolved before a message or action is accepted.

\subsection{Pilot-to-Pilot and Human Communication}

When both parties use Pilot, communication is delivered directly to the
recipient's persona-bound Pilot inbox. Structured requests may include a
human-readable message and a machine-readable proposal. For example, ``ask Alex
to collect the dry cleaning'' becomes a recipient-addressed task proposal, not
an immediate modification of Alex's private list. Alex or Alex's governed Pilot
must accept it before it becomes an assigned task.

When the recipient does not use Pilot, Pilot may prepare a message for an
authorized email, SMS, WhatsApp or equivalent adapter. It must resolve the exact
contact privately, display the exact channel, recipient and content, obtain any
required approval, execute through an officially permitted interface, verify a
provider receipt and only then report delivery. A useful handoff may include an
invitation to Pilot, but joining Pilot is never required merely to read an
ordinary message sent through an existing channel.

This provides the intended adoption loop:

\begin{enumerate}
  \item a user asks Pilot to coordinate something useful with another person;
  \item the recipient receives a clear request through Pilot or an existing
        authorized channel;
  \item a Pilot recipient may accept, decline, delegate, schedule or ask their
        own agent to handle the request;
  \item a non-user can respond normally and may optionally activate Pilot to
        obtain the richer agent-to-agent workflow; and
  \item both sides receive only the state and receipts they are authorized to
        see.
\end{enumerate}

\subsection{Voice Notes, Push-to-Talk and Calls}

Push-to-talk is retained as a first-class communication mode rather than a
decorative voice button. Holding the control requests a short-lived floor lock;
successful ownership starts bounded audio capture; the lock is renewed only
while the sender remains active; release finalizes and sends the voice payload;
and expiry frees the channel after interruption or disconnection. The visible
interface must distinguish requesting, recording, sending, delivered, failed
and waiting-for-floor states.

Voice notes, push-to-talk frames and call signalling use the same persona,
conversation and space boundary as text. Calls may use WebRTC for media
transport, but the current demonstration frame-encryption and placeholder
insertable-stream implementation are not sufficient for a production claim of
end-to-end encrypted calling. Production calls require audited key agreement,
authenticated participants, key rotation, replay protection, device revocation,
clear recording indicators and tested behaviour across supported browsers and
native clients.

Candidate uses include personal and family conversation, household channels,
field teams, external workspaces, Boardroom discussion, connectivity-loss mode
and Pilot Guardian. Emergency permissions must remain distinct from ordinary
contact permissions.

\subsection{Transport Selection and Wave Operation}

The word ``Wave'' describes a transport-independent identity and delivery
model. It does not imply that an ordinary telephone or laptop can transmit over
special long-range radio frequencies without compatible hardware. Pilot will
select a route according to availability, policy, recipient reachability,
latency, size and sensitivity:

\begin{enumerate}
  \item direct local-network delivery when both authorized nodes are reachable;
  \item the user's private VPN or customer-selected secure tunnel;
  \item normal encrypted internet relay where permitted;
  \item a paired GlyphNet radio bridge and store-and-forward mesh where
        compatible hardware is present; or
  \item an authorized external messaging provider for a non-Pilot recipient.
\end{enumerate}

The existing radio-node design supports serial-connected hardware, regional
band profiles, bounded frame sizes, fragmentation, queues, reconnects and a
bridge based on an ESP32-class or Raspberry Pi-class device with a compatible
radio such as an SX1262 LoRa module. Until physical radio hardware is installed
and qualified, the product operates through normal local or internet networking.

Radio support will initially remain an advanced capability flag. Production
activation requires complete inbound payload reconstruction, authenticated
bridge pairing, real driver tests, encryption that is independent of the radio,
message expiry, duplicate suppression, congestion handling, regional duty-cycle
compliance and end-to-end field tests with the internet disabled. BLE and
Wi-Fi Direct routing remain later transport adapters rather than claimed
capabilities of the first release.

\subsection{GlyphChain as a Selective Proof Rail}

GlyphChain is useful immediately when restricted to proofs of governed activity.
For a consequential action, Pilot creates a local authorization and execution
record. Once the result is known, it may commit a compact, tamper-evident proof
containing:

\begin{itemize}
  \item opaque request, action and receipt identifiers;
  \item hashes of the approved scope and relevant result;
  \item pseudonymous actor, owner, device and executing-agent references;
  \item policy and capability-lease versions;
  \item timestamps, expiry, idempotency and replay information;
  \item the adapter or transport class used;
  \item success, failure, cancellation or recipient-acceptance state; and
  \item a reference to a mother-private provider receipt where applicable.
\end{itemize}

The following material must never be written directly to the chain:

\begin{itemize}
  \item message bodies, voice recordings or call media;
  \item documents, photographs, attachments or extracted private text;
  \item names, telephone numbers, email addresses or address-book contents;
  \item passwords, tokens, cookies, private keys, payment-card data or account
        identifiers;
  \item private memory, calendar descriptions, medical information or raw
        location history; and
  \item plaintext task, purchase, booking or Boardroom contents.
\end{itemize}

The proof rail records that a precisely scoped event occurred; the encrypted
source record remains on the owner-controlled mother brain. Verification can
recompute a hash from an authorized local record without disclosing that record
to the chain. Retention, deletion and legal-hold policy apply to the private
record separately from the immutable proof. Where erasure is required, the
remaining chain commitment must be deliberately non-identifying and unusable to
reconstruct the deleted content.

\subsection{PHO and the Dormant Economic Layer}

The PHO and PhotonPay work is retained because its schemas may later support
machine-to-machine settlement, offline community exchange, metered capabilities
or incentives. It is not activated in the initial Pilot build. The following
must remain unavailable to ordinary users and agents:

\begin{itemize}
  \item minting, transfers and balance mutation through development routes;
  \item demonstration accounts, balances and transaction signatures;
  \item token purchasing, exchange, staking, bonding or investment language;
  \item offline credit that could be mistaken for settled funds; and
  \item any interface implying that a simulated chain commit constitutes a real
        payment.
\end{itemize}

Future activation is a separate programme requiring audited cryptography and
key custody, deterministic settlement, conflict and double-spend handling,
economic and governance review, recovery procedures, operational monitoring,
consumer disclosures, jurisdiction-specific legal and regulatory review, and
an explicit product decision. Messaging and the proof rail must continue to
operate when the complete economic layer is disabled.

\subsection{Production Security Requirements}

Before public release, the reused implementation must be brought under the same
identity and capability model as the AION Device Fabric. The minimum security
work is:

\begin{enumerate}
  \item replace fixed development tokens with signed, persona-bound device
        possession and expiring capability leases;
  \item replace the static Wave directory with private, recoverable and
        revocable Pilot identities;
  \item encrypt message and attachment storage at rest with keys unavailable to
        shared television and peripheral nodes;
  \item use authenticated end-to-end encryption for message payloads and audited
        key establishment rather than demonstration QKD shims;
  \item authenticate every WebSocket, relay and radio-bridge session and bind it
        to an exact identity, space and permitted operation;
  \item enforce message-size, attachment, rate, replay, expiry and duplicate
        limits before persistence or fan-out;
  \item complete inbound radio payload reconstruction and verify integrity
        before delivering an application envelope;
  \item remove mock signatures, placeholder settlement and development wallet
        authority from production builds;
  \item complete and audit call-media encryption before presenting calls as
        end-to-end encrypted;
  \item provide device loss, session revocation, account recovery, export,
        deletion, blocking and abuse-reporting controls; and
  \item resolve the standalone GlyphNet Browser build dependency and reproduce
        signed builds from a locked dependency manifest.
\end{enumerate}

Ledger proofs are not a substitute for these controls. A valid proof can show
that a vulnerable system performed an action; it cannot make the action safe.

\subsection{First Build Scope}

The initial integration will concentrate on the useful, already demonstrated
communications path:

\begin{enumerate}
  \item define and version the normalized GlyphNet/Pilot envelope;
  \item map production Pilot personas, trusted devices and spaces onto GlyphNet
        sender, recipient and channel identities;
  \item place text conversations, agent requests, delegated tasks, voice notes
        and push-to-talk in one Pilot Inbox service;
  \item expose the same inbox through the mobile Personal, Workspace and
        Boardroom modes without merging their authorization boundaries;
  \item persist encrypted conversation state on the owner-selected mother brain;
  \item implement direct local/VPN delivery and a production relay abstraction;
  \item connect approved non-Pilot email or messaging handoffs through the
        existing governed service-adapter pattern;
  \item emit private delivery receipts and selective GlyphChain proof receipts
        for consequential actions;
  \item keep radio transport detectable but disabled unless a qualified bridge
        is explicitly enrolled; and
  \item exclude every PHO, PhotonPay and wallet route from the normal application
        manifest and consumer interface.
\end{enumerate}

The first acceptance demonstration is:

\begin{quote}
A user asks Pilot to send another named Pilot user a task. Pilot resolves the
recipient privately, shows the exact proposal, obtains the required approval,
delivers an encrypted structured request, allows the recipient to accept or
decline it, synchronizes the resulting task state, records delivery and
acceptance receipts, and commits only non-identifying proof hashes to
GlyphChain. The same flow fails closed when identity, authority or delivery
cannot be verified.
\end{quote}

\subsection{Deferred Work Register}

The following capabilities are intentionally recorded for later implementation
or qualification rather than lost:

\begin{itemize}
  \item native iOS and Android background delivery and push notifications;
  \item audited end-to-end encrypted audio and video calling;
  \item encrypted multi-device history synchronization and recovery;
  \item contact discovery that does not upload a readable address book;
  \item message reactions, group administration, moderation and abuse controls;
  \item complete BLE and Wi-Fi Direct transports;
  \item physical LoRa/radio bridge qualification, offline routing and regional
        certification;
  \item delayed, disconnected and intermittently connected agent-to-agent work;
  \item Guardian-grade emergency push-to-talk and connectivity fallback;
  \item optional Wave-address federation between independently hosted Pilot
        mothers;
  \item production consensus, validator operations and proof availability;
  \item PHO, PhotonPay, TESS, wrapped GLYPH, bonds and any other economic
        activation programme; and
  \item independent security, privacy, regulatory and adversarial review.
\end{itemize}

\subsection{Acceptance and Truthfulness Criteria}

The communication layer is not complete merely because a message appears in a
local user interface. Completion requires independently testable evidence that:

\begin{itemize}
  \item the correct persona and space authorized the exact communication;
  \item an unauthorized household, workspace or Boardroom identity cannot read
        or alter it;
  \item reconnect and retry do not duplicate a task, message or external action;
  \item the recipient or provider acknowledged delivery;
  \item acceptance, rejection, cancellation and failure propagate correctly;
  \item private content remains absent from GlyphChain and diagnostic logs;
  \item the useful messaging product works with PHO and radio features disabled;
  \item the local/VPN path remains useful without a paid AI provider;
  \item offline radio claims are made only after a real enrolled bridge completes
        an internet-disconnected field test; and
  \item the interface never labels prepared, queued or simulated work as sent,
        paid, settled, accepted or completed.
\end{itemize}

This boundary allows Pilot to capitalize on the earlier AION and GlyphNet work
immediately while keeping experimental economics and unfinished transports out
of the consumer trust boundary. It also creates the communications foundation
for the intended wider product: one privately owned intelligence that can
coordinate a person's home, personal life, professional work and organisations
through a simple, familiar inbox.
